\Needspace{11\baselineskip}
\begin{center}
\textit{DIGRESSION CONTAINING A TREATMENT OF TERNARY FORMS.}
\end{center}

\Needspace{8\baselineskip}
\Article{266}
So far we have restricted our investigation to functions of the second degree involving \textit{two} indeterminates, and there has been no need to give them a special name. But we can clearly regard this subject as a very special part of the most general investigation \textit{of rational integral algebraic h\,o\,m\,o\,g\,e\,n\,e\,o\,u\,s functions of several indeterminates and various degrees}, and conveniently distinguish these functions by their degree as \textit{forms of the second, third, fourth degree, etc.}, and by the number of indeterminates as \textit{binary, ternary, quaternary forms}, etc. Thus the forms that we have so far called simply forms will be called \textit{binary forms of the second degree}; functions such as
\[
Axx+2Bxy+Cyy+2Dxz+2Eyz+Fzz
\]
where $A$, $B$, $C$, $D$, $E$, $F$ denote given integers) will be called \textit{ternary forms of the second degree}, and so on. This section is primarily devoted to binary forms of the second degree alone; but several truths concerning them, and very beautiful ones, still remain whose true source must be sought in the theory of ternary forms of the second degree. So we insert a short digression on that theory here, in which we will present, from its first elements, what is needed to complete the theory of binary forms. We \SourceLeaf{305}hope that mathematicians will find this more welcome than if we either suppressed those truths or derived them by less natural methods. We must reserve a fuller investigation of this very important subject for another occasion, both because its richness would already take us far beyond the limits of this work, and because we hope that it will be enriched by further substantial advances in the future. Quaternary, quinary, etc.\ forms of the second degree, and all forms of higher degrees, are entirely excluded from our purpose here\HistoricalNote{*)}{For this reason, whenever we speak simply of binary or ternary forms from now on, forms \textit{of the second degree} are always meant.}. Let it be enough to draw mathematicians' attention to this vast field, in which they will find abundant material for exercising their powers and enriching higher arithmetic with outstanding advances.

\Needspace{8\baselineskip}
\Article{267}
It will greatly help clarity to establish a fixed order among the three indeterminates entering a ternary form, just as for binary forms, so that the \textit{first, second} and \textit{third indeterminate} are distinguished from one another. In arranging the individual terms of a form, we will always follow this order: first comes the term involving the square of the first indeterminate, followed in turn by those involving the square of the second, the square of the third, twice the product of the second and third, twice the product of the first and third, and twice the product of the first and second. Finally, the given integers by which these squares and double products are multiplied will be called, in the same order, the \textit{first, second, third, fourth, fifth, sixth coefficient}. Thus
\[
axx+a'x'x'+a''x''x''+2bx'x''+2b'xx''+2b''xx'
\]
will be a ternary form in the proper order, whose first indeterminate is $x$, second $x'$, third $x''$, first coefficient $a$, etc., fourth $b$, etc.\ But it will greatly shorten our notation if we do not always need to denote the indeterminates of a ternary form by special letters. So, when we disregard the indeterminates, we will also denote the same form by
\[
\begin{pmatrix}
a, & a', & a'' \\
b, & b', & b''
\end{pmatrix}
\]
\SourceLeaf{306}

Setting
\[
\begin{array}{lll}
bb-a'a''=A, & b'b'-aa''=A', & b''b''-aa'=A'' \\
ab-b'b''=B, & a'b'-bb''=B', & a''b''-bb'=B''
\end{array}
\]
gives another form
\[
\begin{pmatrix}
A, & A', & A'' \\
B, & B', & B''
\end{pmatrix}\ldots F
\]
which we will call the \textit{adjoint of the form}
\[
\begin{pmatrix}
a, & a', & a'' \\
b, & b', & b''
\end{pmatrix}\ldots f
\]
We then find, denoting for brevity the number
\[
abb+a'b'b'+a''b''b''-aa'a''-2bb'b''\quad\text{by}\quad D
\]
\[
\begin{array}{lll}
BB-A'A''=aD, & B'B'-AA''=a'D, & B''B''-AA'=a''D \\
AB-B'B''=bD, & A'B'-BB''=b'D, & A''B''-BB'=b''D
\end{array}
\]
This shows that the adjoint of the form $F$ is the form
\[
\begin{pmatrix}
aD, & a'D, & a''D \\
bD, & b'D, & b''D
\end{pmatrix}
\]
The number $D$, on whose nature the properties of the ternary form $f$ chiefly depend, will be called the \textit{determinant} of this form. In this way the determinant of the form $F$ becomes $=DD$, or equal to the square of the determinant of the form $f$, of which it is the adjoint.

Thus, \textit{for example}, the ternary form $\left(\begin{smallmatrix}29,\allowbreak &13,\allowbreak &9\\7,\allowbreak &-1,\allowbreak &14\end{smallmatrix}\right)$ has the adjoint $\left(\begin{smallmatrix}-68,\allowbreak &-260,\allowbreak &-181\\217,\allowbreak &-111,\allowbreak &133\end{smallmatrix}\right)$, and the determinant of each $=1$.

Ternary forms of determinant $0$ will be entirely excluded from the following investigation, since, as will be shown in the theory of ternary forms to be presented more fully on another occasion, they are ternary only \textit{in appearance} and are in fact equivalent to binary forms.

\Needspace{8\baselineskip}
\Article{268}
If a ternary form $f$ of determinant $D$, whose indeterminates are $x$, $x'$, $x''$ (the first $=x$, etc.), is transformed into a ternary form $g$ of determinant $E$, whose indeterminates are $y$, $y'$, $y''$, by a substitution of the form
\SourceLeaf{307}
\[
\begin{aligned}
x&=\alpha y+\beta y'+\gamma y'' \\
x'&=\alpha'y+\beta'y'+\gamma'y'' \\
x''&=\alpha''y+\beta''y'+\gamma''y''
\end{aligned}
\]
where all nine coefficients $\alpha$, $\beta$, etc.\ are assumed to be integers, then, for brevity, omitting the indeterminates, we will simply say that $f$ passes into $g$ by the substitution
$(S)$
\[
\begin{array}{ccc}
\alpha, & \beta, & \gamma \\
\alpha', & \beta', & \gamma' \\
\alpha'', & \beta'', & \gamma''
\end{array}
\]
and that $f$ \textit{contains} $g$, or that $g$ \textit{is contained in $f$}. This assumption immediately gives six equations for the six coefficients in $g$, which we need not write down. An easy calculation then gives the following conclusions:

I. Denoting for brevity the number
\[
\alpha\beta'\gamma''+\beta\gamma'\alpha''+\gamma\alpha'\beta''
-\gamma\beta'\alpha''-\alpha\gamma'\beta''-\beta\alpha'\gamma''\quad\text{by }k
\]
we find, after the necessary reductions, $E=kkD$. Thus $D$ divides $E$, and the quotient is a square. So the number $k$ plays a role in transformations of ternary forms similar to that of the number $\alpha\delta-\beta\gamma$ in Article 157 for transformations of binary forms: it is a square root of the quotient of the determinants. We might therefore conjecture that the difference in the sign of $k$ also establishes here an essential distinction between proper and improper transformations and containments. But on looking more closely, it is clear. $f$ also passes into $g$ by the substitution
\[
\begin{array}{ccc}
-\alpha, & -\beta, & -\gamma \\
-\alpha', & -\beta', & -\gamma' \\
-\alpha'', & -\beta'', & -\gamma''
\end{array}
\]
and in the value of $k$, replacing $\alpha$ by $-\alpha$, $\beta$ by $-\beta$, etc.\ gives $-k$. So this substitution would be dissimilar to the substitution $S$, and every ternary form containing another in one way would also contain it in the other. Such a \SourceLeaf{308}distinction, being of no use for ternary forms, will therefore be entirely ruled out here.

II. Denoting by $F$, $G$ the adjoints of $f$, $g$, respectively, the coefficients in $F$ are determined by the coefficients in $f$, and the coefficients in $G$ by the values of the coefficients of the form $g$, known from the equations supplied by the substitution $S$. Expressing the coefficients of the form $f$ by letters and comparing the values of the coefficients of the forms $F$, $G$, we readily confirm that $F$ contains the form $G$ and is transformed into it by the substitution
$(S')$
\[
\begin{array}{ccc}
\beta'\gamma''-\beta''\gamma', & \gamma'\alpha''-\gamma''\alpha', & \alpha'\beta''-\alpha''\beta' \\
\beta''\gamma-\beta\gamma'', & \gamma''\alpha-\gamma\alpha'', & \alpha''\beta-\alpha\beta'' \\
\beta\gamma'-\beta'\gamma, & \gamma\alpha'-\gamma'\alpha, & \alpha\beta'-\alpha'\beta
\end{array}
\]
We omit the calculation itself, which presents no difficulties.

III. The form $g$ is, by the substitution $(S'')$
\[
\begin{array}{ccc}
\beta'\gamma''-\beta''\gamma', & \beta''\gamma-\beta\gamma'', & \beta\gamma'-\beta'\gamma \\
\gamma'\alpha''-\gamma''\alpha', & \gamma''\alpha-\gamma\alpha'', & \gamma\alpha'-\gamma'\alpha \\
\alpha'\beta''-\alpha''\beta', & \alpha''\beta-\alpha\beta'', & \alpha\beta'-\alpha'\beta
\end{array}
\]
clearly transformed into the same form into which $f$ passes by
\[
\begin{array}{ccc}
k, & 0, & 0 \\
0, & k, & 0 \\
0, & 0, & k
\end{array}
\]
that is, into the form obtained by multiplying each coefficient of the form $f$ by $kk$. We will denote this form by $f'$.

IV. In exactly the same way it is proved that the form $G$, by the substitution
$(S''')$
\[
\begin{array}{ccc}
\alpha, & \alpha', & \alpha'' \\
\beta, & \beta', & \beta'' \\
\gamma, & \gamma', & \gamma''
\end{array}
\]
\SourceLeaf{309}
passes into the form obtained from $F$ by multiplying each coefficient by $kk$. We will denote this form by $F'$.

We will say that the substitution $S'''$ arises \textit{by transposing} the substitution $S$. Clearly, $S$ then arises in turn by transposing the substitution $S'''$, and $S'$, $S''$ each by transposing the other. --- The substitution $S'$ can conveniently be called the \textit{adjoint} of the substitution $S$, so the adjoint of the substitution $S'''$ will be $S''$.

\Needspace{8\baselineskip}
\Article{269}
If not only the form $f$ contains $g$, but also the latter contains the former, the forms $f$, $g$ will be called \textit{equivalent}. In this case not only will $D$ divide $E$, but $E$ will also divide $D$, from which we easily conclude that $D=E$ must hold. Conversely, if a form $f$ contains a form $g$ of the same determinant, these two forms will be equivalent. For (using the same notation as in the preceding article and excluding the case where $D=0$) we will have $k=\pm1$, and so the form $f'$. into which $g$ passes by the substitution $S''$, is identical to $f$, or $f$ is contained in $g$. It is also clear that in this case the forms $F$, $G$, adjoints of $f$, $g$, will be equivalent to each other, and the latter will pass into the former by the substitution $S'''$. Finally, conversely, if the forms $F$, $G$ are \textit{assumed} equivalent, and the former passes into the latter by the substitution $T$, the forms $f$, $g$ will also be equivalent: $f$ will pass into $g$ by the adjoint of the substitution $T$, and $g$ into $f$ by the transpose of the substitution $T$. For by these two substitutions, respectively, the adjoint of $F$ passes into the adjoint of $G$, and the latter into the former. But these two forms are obtained from $f$, $g$ by multiplying each coefficient by $D$, so we readily conclude that, by the same substitutions, $f$ passes into $g$. and $g$ into $f$, respectively.

\Needspace{8\baselineskip}
\Article{270}
If a ternary form $f$ contains a ternary form $f'$, and the latter a form $f''$: then $f$ will also contain $f''$. For it is very easy to see that if the following transformations hold:
\[
\begin{array}{c|c}
f\text{ into }f'\text{ by substitution} & f'\text{ into }f''\text{ by substitution} \\[2pt]
\begin{array}{ccc}
\alpha, & \beta, & \gamma \\
\alpha', & \beta', & \gamma' \\
\alpha'', & \beta'', & \gamma''
\end{array}
&
\begin{array}{ccc}
\delta, & \varepsilon, & \zeta \\
\delta', & \varepsilon', & \zeta' \\
\delta'', & \varepsilon'', & \zeta''
\end{array}
\end{array}
\]
\SourceLeaf{310}
$f$ will be transformed by the substitution
\[
\begin{array}{ccc}
\alpha\delta+\beta\delta'+\gamma\delta'', & \alpha\varepsilon+\beta\varepsilon'+\gamma\varepsilon'', & \alpha\zeta+\beta\zeta'+\gamma\zeta'' \\
\alpha'\delta+\beta'\delta'+\gamma'\delta'', & \alpha'\varepsilon+\beta'\varepsilon'+\gamma'\varepsilon'', & \alpha'\zeta+\beta'\zeta'+\gamma'\zeta'' \\
\alpha''\delta+\beta''\delta'+\gamma''\delta'', & \alpha''\varepsilon+\beta''\varepsilon'+\gamma''\varepsilon'', & \alpha''\zeta+\beta''\zeta'+\gamma''\zeta''
\end{array}
\]

Thus, in the case where $f$ is equivalent to $f'$, and $f'$ to $f''$, the form $f$ will also be equivalent to the form $f''$. --- It is also immediately clear how these theorems are to be applied to more forms.
